% =============================================================================
% Copyright (c) 2026 Mert Torun, M.Sc. (mtorun0x7cd) <info@mtorun0x7cd.com>
% 
% Permission is hereby granted, free of charge, to any person obtaining a copy
% of this software and associated documentation files (the "Software"), to deal
% in the Software without restriction, including without limitation the rights
% to use, copy, modify, merge, publish, distribute, sublicense, and/or sell
% copies of the Software, and to permit persons to whom the Software is
% furnished to do so, subject to the following conditions:
% 
% The above copyright notice and this permission notice shall be included in
% all copies or substantial portions of the Software.
% 
% THE SOFTWARE IS PROVIDED "AS IS", WITHOUT WARRANTY OF ANY KIND, EXPRESS OR
% IMPLIED, INCLUDING BUT NOT LIMITED TO THE WARRANTIES OF MERCHANTABILITY,
% FITNESS FOR A PARTICULAR PURPOSE AND NONINFRINGEMENT. IN NO EVENT SHALL THE
% AUTHORS OR COPYRIGHT HOLDERS BE LIABLE FOR ANY CLAIM, DAMAGES OR OTHER
% LIABILITY, WHETHER IN AN ACTION OF CONTRACT, TORT OR OTHERWISE, ARISING FROM,
% OUT OF OR IN CONNECTION WITH THE SOFTWARE OR THE USE OR OTHER DEALINGS IN THE
% SOFTWARE.
% =============================================================================

\chapter{Source Code Listings}
\label{app:code}

This appendix provides a high-level structural map and documents key low-level code implementations of the \textit{OneWare Container Extension}. The complete codebase is tracked in the repository \url{https://github.com/FEntwumS/FEntwumS.ContainerExtension} on the \texttt{main} branch.

\section{High-Level Architectural Mapping}
To aid in navigating the project source code, Table~\ref{tab:directory_map} maps the directory structure to the key architectural classes and their designated roles within the extension.

\begin{table}[H]
    \caption{Mapping of Directory Structure to Key Architectural Classes}
    \label{tab:directory_map}
    \centering
    \begin{tabularx}{\textwidth}{@{} p{4.6cm} p{4.2cm} X @{}}
    \toprule
    \textbf{Directory/Path} & \textbf{Key Class} & \textbf{Architectural Role} \\
    \midrule
    \path{src/ContainerExtension/} & \texttt{DockerExecutionStrategy} & Core coordinator, dynamic platform resolution, and lifecycle manager. \\
    \path{src/ContainerExtension/} & \texttt{ContainerTelemetry} & Thread-safe, low-allocation structured JSON telemetry logger. \\
    \path{src/ContainerExtension/.../Docker/} & \texttt{DockerCommandBuilder} & Builds the shell-less argument vector, host-to-container path mappings, and volume/security configuration. \\
    \path{src/ContainerExtension/Registry/} & \texttt{RegistryClient} & Handles OCI registry HTTP communication, authentication, and registry tag enumeration (SSRF-guarded). \\
    \bottomrule
    \end{tabularx}
\end{table}

\lstset{
  style=csharp-appendix,
  literate={—}{---}1 {⚠️}{[WARN]}1
}

\section{Windows Named Pipe Privilege Escalation Override}
Listing~\ref{lst:secure_named_pipe_get_handler} illustrates how .NET Reflection is utilized in the C\# client to target the non-public \texttt{\_streamOpener} field of \texttt{ManagedHandler} and overwrite it with our custom stream opener.

\begin{lstlisting}[language=CSharp, caption={[Reflection-based Handler Interception in SecureNamedPipeCredentials]{Reflection-based Handler Interception in SecureNamedPipeCredentials.GetHandler}}, label={lst:secure_named_pipe_get_handler}]
public override System.Net.Http.HttpMessageHandler GetHandler(System.Net.Http.HttpMessageHandler innerHandler)
{
    if (string.Equals(innerHandler.GetType().FullName, "Microsoft.Net.Http.Client.ManagedHandler", StringComparison.Ordinal))
    {
        var field = innerHandler.GetType().GetField("_streamOpener", 
            System.Reflection.BindingFlags.NonPublic | System.Reflection.BindingFlags.Instance);
        if (field != null)
        {
            var delegateType = field.FieldType;
            var method = typeof(SecureNamedPipeCredentials).GetMethod(nameof(SecureStreamOpenerAsync), 
                System.Reflection.BindingFlags.NonPublic | System.Reflection.BindingFlags.Instance);
            if (method != null)
            {
                var d = System.Delegate.CreateDelegate(delegateType, this, method);
                field.SetValue(innerHandler, d);
            }
        }
    }
    return innerHandler;
}
\end{lstlisting}

Listing~\ref{lst:secure_stream_opener_async} details the custom socket stream opener that establishes the named pipe client stream using identification-level impersonation to block local privilege escalation.

\begin{lstlisting}[language=CSharp, caption={Custom SecureStreamOpenerAsync with Identification Impersonation}, label={lst:secure_stream_opener_async}]
private async System.Threading.Tasks.Task<System.IO.Stream> SecureStreamOpenerAsync(string host, int port, System.Threading.CancellationToken token)
{
    var pipeName = _endpoint.LocalPath;
    var serverName = ".";
    if (pipeName.StartsWith(@"\\", StringComparison.Ordinal))
    {
        var parts = pipeName.Split('\\', StringSplitOptions.RemoveEmptyEntries);
        if (parts.Length >= 3)
        {
            serverName = parts[0];
            pipeName = parts[parts.Length - 1];
        }
    }
    else
    {
        if (pipeName.StartsWith("pipe/", StringComparison.OrdinalIgnoreCase))
            pipeName = pipeName[5..];
        else if (pipeName.StartsWith("/pipe/", StringComparison.OrdinalIgnoreCase))
            pipeName = pipeName[6..];
    }

    var pipe = new System.IO.Pipes.NamedPipeClientStream(
        serverName,
        pipeName,
        System.IO.Pipes.PipeDirection.InOut,
        System.IO.Pipes.PipeOptions.Asynchronous,
        System.Security.Principal.TokenImpersonationLevel.Identification);

    try
    {
        await pipe.ConnectAsync(token).ConfigureAwait(false);
        return pipe;
    }
    catch
    {
        await pipe.DisposeAsync().ConfigureAwait(false);
        throw;
    }
}
\end{lstlisting}

\section{Path Resolution Security Verification}
Listing~\ref{lst:resolve_canonical_internal} displays the recursive segment-by-segment canonicalization logic used to resolve symbolic links and check path safety boundaries. The excerpt is condensed: the production detector branches on \texttt{Directory.Exists} and \texttt{File.Exists} to cover both file and directory reparse points, whereas the listing shows only the directory branch for brevity.

\begin{lstlisting}[language=CSharp, caption={Recursive Symbolic Link Resolution (ResolveCanonicalInternal, condensed)}, label={lst:resolve_canonical_internal}]
string ResolveCanonicalInternal(string currentPath, int depth)
{
    if (depth > 40)
        throw new ContainerExtension.DockerExecutionException("Too many levels of symbolic links.");

    string absolutePath = currentPath;
    if (!Path.IsPathRooted(absolutePath))
        absolutePath = Path.Combine(Directory.GetCurrentDirectory(), absolutePath);

    string root = Path.GetPathRoot(absolutePath) ?? (OperatingSystem.IsWindows() ? "C:\\" : "/");
    string remainder = absolutePath.Substring(root.Length);
    var components = remainder.Split(new char[] { '/', '\\' }, StringSplitOptions.RemoveEmptyEntries);
    string current = root;

    foreach (var component in components)
    {
        if (string.Equals(component, ".", StringComparison.Ordinal)) continue;
        if (string.Equals(component, "..", StringComparison.Ordinal))
        {
            var parent = Path.GetDirectoryName(current);
            current = parent ?? root;
            continue;
        }

        string next = Path.Combine(current, component);
        bool isSymlink = false;
        string? target = null;

        try
        {
            var info = new DirectoryInfo(next);
            if (info.Exists && info.Attributes.HasFlag(FileAttributes.ReparsePoint))
            {
                target = info.LinkTarget;
                isSymlink = !string.IsNullOrEmpty(target);
            }
        }
        catch { /* fallback check logic */ }

        if (isSymlink && target != null)
        {
            string canonicalSymlink = Path.GetFullPath(next);
            if (!seenSymlinks.Add(canonicalSymlink))
                throw new ContainerExtension.DockerExecutionException($"Circular symbolic link detected: '{next}'");

            try
            {
                current = Path.IsPathRooted(target) 
                    ? ResolveCanonicalInternal(target, depth + 1)
                    : ResolveCanonicalInternal(Path.Combine(current, target), depth + 1);
            }
            finally { seenSymlinks.Remove(canonicalSymlink); }
        }
        else
        {
            current = next;
        }
    }
    return Path.GetFullPath(current);
}
\end{lstlisting}

\section{Asynchronous Error Telemetry Channel Processing}
Listing~\ref{lst:process_error_channel} demonstrates the asynchronous background processing loop using .NET \texttt{System.\allowbreak Threading.\allowbreak Channels} to serialize error logging actions asynchronously.

\begin{lstlisting}[language=CSharp, caption={Asynchronous Channel Telemetry consumer loop}, label={lst:process_error_channel}]
private static System.Threading.Channels.Channel<TelemetryErrorEntry> ErrorChannel =
    System.Threading.Channels.Channel.CreateUnbounded<TelemetryErrorEntry>(new System.Threading.Channels.UnboundedChannelOptions
    {
        SingleReader = true,
        SingleWriter = false
    });

static ContainerTelemetry()
{
    _ = Task.Run(ProcessErrorChannelAsync);
}

private static async Task ProcessErrorChannelAsync()
{
    var reader = ErrorChannel.Reader;
    await foreach (var entry in reader.ReadAllAsync().ConfigureAwait(false))
    {
        WriteErrorEntryToDisk(entry);
    }
}
\end{lstlisting}
